\section{Introduction}\label{sec:introduction}
In a sequence of independent fair signs, a long constant run is rare for a
simple reason: each additional sign must agree with the preceding one.
A left-maximal run of length at least $L$ starts with probability $2^{-L}$,
so $N$ potential starts have natural intensity $N2^{-L}$. The transition
between many runs and no runs occurs at $L=\log_2N+O(1)$; recording starts
rather than all constant windows removes the local clustering within a run
\cite{ArratiaGoldsteinGordon1990,ArratiaGordonWaterman1990,Godbole1991}.

Logarithmic blocks also underlie the original Erd\H{o}s--R\'enyi law
of large numbers \cite{ErdosRenyi1970}; the distributional and marked
run laws used here are the later results cited above.

Now impose complete multiplicativity. Choose the signs $f(p)$ independently
at the primes and set
\begin{equation}\label{eq:model}
 f(1)=1,\qquad f(n)=\prod_p f(p)^{v_p(n)},\qquad
 \PP(f(p)=1)=\PP(f(p)=-1)=\tfrac12.
\end{equation}
In particular $f(n^2)=1$. This is the completely multiplicative, sometimes
called \emph{extended Rademacher}, model, not the squarefree-supported model. Random
multiplicative models go back to Wintner \cite{Wintner1944}; the
completely multiplicative convention here is fixed by \eqref{eq:model}.
The identity $f(2n)=f(2)f(n)$ persists at every scale. Distant blocks therefore cannot simply be treated as independent: the
effect of these identities must be counted. Nevertheless, the critical run
statistics have the independent-sign Poisson form.

There is a second, quite different mechanism at the left boundary. Requiring
$f(1),\ldots,f(L)$ to be equal is exactly the requirement that every prime
sign up to $L$ be positive. Its probability is $2^{-\pi(L)}$, not $2^{-L}$.
Thus a sufficiently exceptional run can be produced either by many potential
starts in the bulk or by a single arithmetically cheaper event at the boundary.
Our results quantify both mechanisms and their competition.

The same approach also distinguishes many prescribed words at once. Two
sources of dependence must be separated: words can overlap because a
suffix matches a prefix, and distant occurrences can interact because a
product of their integer arguments is a square. The first is a property
of the dictionary; the second is a property of the arithmetic sequence.
Keeping them separate is the organizing principle of the paper.

\subsection{Main results}
For $x\ge2$ and $L\ge1$, let
\begin{equation}\label{eq:start}
 J_{x,L}=\1_{\{f(x-1)=-f(x),\ f(x)=\cdots=f(x+L-1)\}}.
\end{equation}
Write $I_N=[N,2N)\cap\Z$, $Z_{N,L}=\sum_{x\in I_N}J_{x,L}$ and
$\lambda_N=N2^{-L}$. On a start, let $\ell_x$ be the maximal rightward run
length. Every such length is almost surely finite. Define
\[
 K_{x,e}=\1_{\{J_{x,L}=1,\ \ell_x=L+e\}},\qquad
 K^s_{x,e}=K_{x,e}\1_{\{f(x)=s\}},\quad e\in\NN,\ s\in\{\pm1\}.
\]
Let $\nu_{\rm geo}\{e\}=2^{-e-1}$ and
$\nu_\pm=(\delta_{-1}+\delta_{+1})/2$.
A \emph{word} is a prescribed finite sign sequence, a \emph{dictionary}
is a set of such words, and \emph{pattern} is the broader term for the
corresponding occurrence observable. All logarithms without an indicated
base are natural. All total-variation distances below use the convention
$\dTV(\mu,\nu)=\sup_A|\mu(A)-\nu(A)|$.

\begin{theorem}[Critical sign-resolved run field]\label{thm:poisson-main}
Fix $C>0$. Uniformly for integer lengths $|L-\log_2N|\le C$,
\begin{equation}\label{eq:poisson-main}
 \dTV\left(\Law((K^s_{x,e})_{x,e,s}),
       \bigotimes_{x,e,s}\Pois(2^{-L-e-2})\right)
 \le \exp\left[-\left(\frac1{\sqrt2}+o_C(1)\right)
                    \sqrt{\log N\log\log N}\right].
\end{equation}
Here $(x,e,s)$ ranges over $I_N\times\NN\times\{\pm1\}$, and the
comparison is on the common countable lattice. In particular the same
upper bound applies to $\dTV(\Law(Z_{N,L}),\Pois(\lambda_N))$.
\end{theorem}
Since $L$ is integral, changing it by one divides $\lambda_N$ by two;
$\lambda_N$ need not converge. This is why the approximation retains
its moving parameter. For continuous positions, the consequence is a weak limit. Along a
subsequence with $\lambda_N\to\lambda\in(0,\infty)$,
\[
 \sum_{x\in I_N}\sum_{e\ge0}\sum_{s\in\{\pm1\}}K^s_{x,e}\delta_{(x/N,e,s)}
 \ \Longrightarrow\ {\rm PPP}(\lambda\,dt\otimes\nu_{\rm geo}\otimes\nu_\pm)
\]
weakly on finite point configurations on $[1,2]\times\NN\times\{\pm1\}$.
The nonempty configurations on a deterministic spatial grid and those of a
diffuse Poisson process are singular. Consequently the last convergence is
weak, not in total variation. The lattice comparison is stronger in the
appropriate prelimit space: a single coupling controls positions, excesses,
signs, threshold counts and order statistics simultaneously. In the bulk,
positive and negative long runs have equal limiting intensities, despite
the deterministic positivity of every square.

Let $R_M$ be the longest constant run in $(f(1),\ldots,f(M))$.
\begin{theorem}[Longest run and competing rare-event sources]
\label{thm:global-summary}
At critical lengths $|L-\log_2M|\le C$,
\begin{equation}\label{eq:intro-prefix}
 \PP(R_M<L)=\exp(-M2^{-L})+o_C(1).
\end{equation}
Moreover,
\begin{equation}\label{eq:intro-as}
 R_M=\log_2M+O(\log\log M)\qquad\hbox{almost surely}.
\end{equation}
Suppose instead that $L\to\infty$, $L=O(\log M)$, and
$\lambda_M=M2^{-L}\to0$. Then
\begin{equation}\label{eq:intro-two-source}
 \PP(R_M\ge L)=2^{-\pi(L)}+M2^{-L}
                  +o(2^{-\pi(L)}+M2^{-L}).
\end{equation}
Let $X_{M,L}$ be the leftmost start of a length-$L$ constant window contained
in $[1,M]$, on the event $R_M\ge L$. If
$L-\log_2M-\pi(L)\to s\in\mathbb R$, then
\begin{equation}\label{eq:intro-location}
 \Law(X_{M,L}/M\mid R_M\ge L)
 \ \Longrightarrow\
 \frac1{1+2^{-s}}\delta_0+
 \frac{2^{-s}}{1+2^{-s}}\Leb_{[0,1]}.
\end{equation}
The boundary branch is more precise than its atom at macroscopic position
zero: with conditional probability tending to one its start is exactly $1$.
\end{theorem}
The centring in \eqref{eq:intro-location} compares two exact
costs, $L-\log_2M$ and $\pi(L)$, rather than replacing the prime count by a
first-order asymptotic at a precision where that replacement is invalid.
\Cref{thm:two-clock} strengthens the location law to a lattice comparison
with moving source weights, sign and overshoot. Its geometric clock counts
integers in the bulk and prime ranks at the boundary; no phase-convergent
subsequence is needed for that prelimit comparison.

The third result treats many words at once, with their absolute signs
prescribed. Let $\mathcal W_N\subset\{\pm1\}^{B_N}$ be a deterministic
set of $m_N\ge1$ distinct words. Define
\[
 I_{x,w}=\1_{\{f(x-1+j)=w_j\ (0\le j<B_N)\}},\qquad
                     x\in I_N,\ w\in\mathcal W_N.
\]
Write $w\sim_d v$ if the suffix $(w_d,\ldots,w_{B_N-1})$ equals the
prefix $(v_0,\ldots,v_{B_N-d-1})$, and set
\begin{equation}\label{eq:intro-overlap}
 \Omega(\mathcal W_N)=\frac1{m_N}\sum_{w,v\in\mathcal W_N}
           \sum_{d=1}^{B_N-1}2^{-d}\1_{\{w\sim_d v\}}.
\end{equation}
\Needspace{14\baselineskip}
\begin{theorem}[Growing dictionaries of rare words]\label{thm:patterns-summary}
Fix $C>0$ and $0<\eta<1/2$. Suppose
\[
 m_N\le N^{1/2-\eta},\qquad
 |B_N-\log_2(Nm_N)|\le C,\qquad
 \Omega(\mathcal W_N)\longrightarrow0.
\]
Then
\begin{equation}\label{eq:intro-word-field}
 \dTV\left(\Law((I_{x,w})_{x\in I_N,w\in\mathcal W_N}),
       \bigotimes_{x\in I_N,w\in\mathcal W_N}\Pois(2^{-B_N})\right)
                          \longrightarrow0.
\end{equation}
The same conclusion holds after any deterministic restriction of the
site--word index set. The rare-word field is also asymptotically equivalent
in total variation to the field formed from independent fair signs.
\end{theorem}
The total target intensity is $Nm_N2^{-B_N}$, while each word has mean
$N2^{-B_N}$. Thus the comparison retains both the location and identity
of a rare word, even when the dictionary grows. The theorem is proved
quantitatively in \cref{thm:word-field}; \cref{cor:generic-dictionaries}
shows that all but an $O(\log N/\sqrt N)$ fraction of dictionaries of the
stated size have a suitably small overlap weight.
There are also explicit deterministic examples with no overlaps at all:
\cref{cor:marker-dictionaries} uses a classical binary marker construction.
For example, $\lfloor N^{2/5}\rfloor$ prescribed words of length
$(7/5)\log_2N+O(1)$ can be treated simultaneously. In this example the
number of word types grows, but the expected total number of occurrences
stays bounded. These are different parameters, and the theorem controls
the former without discarding any labels.

For a single word counted up to sign, take $\mathcal W=\{a,-a\}$.
The run-start word $a=(-1,+1,\ldots,+1)$ has overlap weight $2^{-L}$.
Dictionaries with nonvanishing weighted overlap are outside the theorem;
short-period words illustrate the possible local clusters.
The compound-Poisson law in \cref{cor:window-clusters} describes the
clusters produced by counting every constant window rather than only
run starts.

The three main results use the same comparison in three forms: the raw
profile controls run fields, the capped profile handles growing families
of mutually exclusive words, and the relative comparison resolves the
competition between the exact boundary event and the bulk.

\subsection{The mechanism and the arithmetic estimate}
Write prime signs as $f(p)=(-1)^{X_p}$ with independent uniform
$X_p\in\mathbb F_2$. A pattern up to sign prescribes $L$ affine linear
conditions on these prime coordinates. For two separated windows, let
$\rho_L(x,y)$ be the dimension of the homogeneous relation space of their
stacked systems. It depends on the windows, not on the prescribed signs.
The affine Fourier identity gives
\[
 \left|\PP(I_x=I_y=1)-2^{-2L}\right|
 \le 2^{-2L}(2^{\rho_L(x,y)}-1).
\]
Thus the useful arithmetic object counts \emph{all nonzero relations},
with their full multiplicity, rather than merely the pairs that carry one.
For prescribed absolute words the relevant kernel omits the two block-parity
conditions. Its dimension increases by at most two. A host count valid
without parity restrictions shows that the same weighted bound survives;
this is the first bridge from relative signs to the dictionary theorem.

\Needspace{4\baselineskip}
The second bridge is probabilistic: at a fixed site the dictionary words are
mutually exclusive, so the chance of a joint occurrence cannot exceed a
single-site probability. Capping the relation weight at this ceiling,
before summing over pairs, is what permits dictionaries of size
$N^{1/2-\eta}$.

The uniform estimate proved in \cref{sec:two-window-anatomy} is, for fixed
$0<\delta<1$ and $B=L+1\asymp\log M$,
\begin{equation}\label{eq:intro-master-profile}
 \sum_{\substack{M^\delta\le x,y<M\\|x-y|>L}}
 (2^{\rho_L(x,y)}-1)
 \ll_\varepsilon M^\varepsilon
 \left(M^{3/2}2^{B/6}+M2^{2B/3}+M^{2/3}2^B\right).
\end{equation}
The fixed logarithmic band and $\delta$ are understood in the implicit
constant. This single geometry includes the dyadic and bounded-ratio
estimates by restriction. At criticality the bound is $M^{5/3+\varepsilon}$.

Two devices make the proof work. Exact rational coincidences are first
removed as a systematic binary code. The remaining relations are resolved
by the graph of prime coordinates larger than $B$. For the hardest terminal
population, determinants force many small odd kernels into the same short
window. Counting the kernels separately is wasteful; a shifted pair-energy
estimate detects the necessary concentration. Finally one conditions on
prime signs below a larger cutoff $Y$. Most starts then have exact
conditional marginal $2^{-L}$, and disjoint prime supports give an exact
conditional dependency graph. Chen--Stein is applied \emph{after} this
arithmetic separation, not as a substitute for it.

\begin{figure}[t]
\centering
\begin{tikzpicture}[node distance=9mm and 9mm]
\node[flowblue,text width=60mm] (a) {Affine word systems (Section~2)\\
\small relative signs and prescribed values};
\node[flowrust,text width=60mm,right=of a] (b) {Rational code and residual graph\\(Section~3)\\
\small CRT, Runge, Pell and kernel energy};
\node[flowblue,text width=60mm,below=of a] (c) {Condition on small primes (Section~4)\\
\small private pivots;\\exact dependency graph};
\node[flowrust,text width=60mm,right=of c] (d) {Poisson comparison (Sections~4--6)\\
\small word fields and signed run marks};
\draw[flowarrow] (a)--(b);
\draw[flowarrow] (b)--(d);
\draw[flowarrow] (a)--(c);
\draw[flowarrow] (c)--(d);
\end{tikzpicture}
\caption{The arithmetic and probabilistic interfaces. The weighted two-window
estimate bounds the separated-pair term. Private pivots supply the exact
conditional graph and the correct marginals; local overlap is treated
separately. Neither interface is dispensable.}\label{fig:proof-flow}
\end{figure}

\subsection{Relation to earlier work}
The independent-sign theory supplies the run-start normalization, the
geometric excess law and the lattice extreme-value comparison
\cite{Foldes1979,GordonSchillingWaterman1986,ArratiaGordonWaterman1990}.
Dependency-graph Poisson approximation is classical
\cite{ArratiaGoldsteinGordon1989,ArratiaGoldsteinGordon1990,ChenRollin2013}.
Poisson approximation for occurrences of prescribed words and the role of
self-overlap are also classical \cite{Godbole1991,GodboleSchaffner1993}.
For several words, Reinert--Schbath \cite{ReinertSchbath1998} give
Poisson-process and compound-Poisson approximations for clumps in Markov
chains; Roquain--Schbath \cite{RoquainSchbath2007} treat overlapping rare
word families. Thus a joint approximation with word types is not itself
a new probabilistic principle.
The contribution here is to verify joint, site-labelled comparisons in a
nonstationary arithmetic sequence with exact multiplicative identities,
uniformly for logarithmic word length and a growing dictionary.

Najnudel proves eventual independence of fixed-length blocks and almost-sure
equidistribution of fixed words in roots-of-unity models \cite{Najnudel2020}.
In the sign model the individual laws can be degenerate at squares, so
this independence does not say that every coordinate is a fair sign. Fixed-pattern and sign-complexity questions also
arise for deterministic multiplicative functions
\cite{ButtkewitzElsholtz2011,MatomakiRadziwillTao2016,Mangerel2024}.
Our growing block length and simultaneous count over moving starts require
uniform weighted control of pairs of windows, beyond a fixed-block
independence statement. The pattern theorem makes explicit which part of that
control is insensitive to the affine right-hand side.

Square-product paucity and Diophantine rigidity are important tools in random
multiplicative functions. Related applications include short-interval central
limit theorems and random Chowla problems
\cite{ChatterjeeSoundararajan2012,SoundararajanXu2023,KlurmanShkredovXu2023,ChinisShala2025}. For fixed polynomial values, de la Bret\`eche, Wang and Xu
use square-product counts and Evertse--Silverman or Pell estimates
\cite{deLaBretecheWangXu2026}. Munsch--Shparlinski--Sun--Xiao
\cite[Theorem~1.3]{MunschShparlinskiSunXiao2026} prove an unconditional
square-product bound, uniform in the translation of one interval, for
each fixed number of factors. The constants depend on that fixed number;
the result does not directly control all relation weights in two moving
windows of logarithmic length. Tao's study of products of consecutive integers
with unusual anatomy provides a particularly close comparison for the
smooth-number and Diophantine mechanisms \cite{Tao2026}.
These works do not provide the growing-window profile
\eqref{eq:intro-master-profile} as an immediate specialization.
The properly normalized short-interval sums of Harper, Soundararajan and Xu
\cite{HarperSoundararajanXu2026} concern a different observable and, in their
stated model, Steinhaus prime values rather than binary run indicators.

We use Runge's method in an explicitly proved split-product form
\cite{Runge1887,Walsh1992}; for fixed-degree component counts, bounded-height square-class localization
and a Pell equation suffice (\cref{lem:split-bounded}). Uniform integral-point
theorems such as Evertse--Silverman~\cite{EvertseSilverman1986} provide
stronger, height-free information, but are not inputs to our proof.
The microscopic analysis uses
the uniform prime-divisor lower bound of Laishram--Shorey
\cite{LaishramShorey2004} and the pointwise result of
Balasubramanian--Shorey \cite{BalasubramanianShorey1993}.
The geometric marked target, its thinning chain and its Gaussian limits are
probabilistic consequences once the arithmetic comparison is established;
they are not presented as new properties of Poisson processes.

\subsection{Reading guide and scope}
Section~\ref{sec:affine-one-window} develops the affine encoding and the
one-window bounds. Section~\ref{sec:two-window-anatomy} proves the final
uniform two-window estimate and its extension to prescribed values.
Sections~\ref{sec:poisson-transfer}--\ref{sec:conditional-laws} extract scalar,
pattern, marked and conditional consequences. Section~\ref{sec:global-prefix}
handles the finite prefix, the boundary and the relative crossover.
A reader primarily interested in probability may take the uniform,
prescribed-value and capped two-window profiles as arithmetic inputs
(\cref{thm:master-profile,prop:absolute-profile,prop:capped-profile}), then
continue with Sections~\ref{sec:poisson-transfer}--\ref{sec:global-prefix}.
The two-window proof includes its Diophantine counting arguments, using
standard ideal and unit theory for the Pell step. Section~\ref{sec:global-prefix}
adds the explicitly stated prime-divisor and smooth-product inputs needed
near the boundary. The companion expands the technical steps called from
the article, rather than giving a second proof of the relation estimate.
The exact prime-cutoff definitions in Section~\ref{sec:poisson-transfer}
are not needed to read the three main statements.

We do not assert optimality of the exponent $5/3$ or of the prime-cutoff
ranges. The full one-factor range for the spatially labelled field remains
unproved. The dictionary-size exponent is likewise a sufficient bound,
not an optimality statement. Records, occupation laws, scale flows and
auxiliary phase removal are developed in the separate process article
\cite{PoulyLatticePoissonFlows2026}; none is a proof input here.
